\documentclass[12pt]{article}
\def\activityid{F06-X-v2}\usepackage[letterpaper,margin=.65in,top=.60in,bottom=.65in]{geometry}
\usepackage[T1]{fontenc}\usepackage[default]{sourcesanspro}
\usepackage{microtype,tikz,array,tabularx,fancyhdr,amsmath}\usepackage[hidelinks]{hyperref}
\usetikzlibrary{arrows.meta,calc}
\setlength{\parindent}{0pt}\setlength{\parskip}{7pt}\setlength{\headheight}{0pt}\setlength{\footskip}{26pt}\setlength{\emergencystretch}{2em}
\pagestyle{fancy}\fancyhf{}\renewcommand{\headrulewidth}{0pt}\renewcommand{\footrulewidth}{.4pt}
\fancyfoot[L]{\scriptsize Bellingham Math Circle / Week 6 / \activityid}\fancyfoot[R]{\scriptsize \thepage}
\definecolor{ink}{gray}{.12}\definecolor{soft}{gray}{.95}\color{ink}
\newcommand{\PageTitle}[2]{{\fontsize{10}{12}\selectfont\bfseries WEEK 6\quad CODE LAB\hfill #1\par}\vspace{3pt}{\fontsize{26}{29}\selectfont\bfseries #2\par}\vspace{2pt}}
\newcommand{\NameLine}{{\small Name: \rule{2.65in}{.4pt}\hfill Date: \rule{1.35in}{.4pt}\par}\vspace{4pt}}
\newcommand{\Task}[2]{\par\vspace{3pt}{\large\bfseries #1}\par #2\par}
\newcommand{\Subhead}[1]{\par\vspace{3pt}{\large\bfseries #1}\par}
\newcommand{\RuleBox}[1]{\par\begingroup\setlength{\fboxsep}{8pt}\colorbox{soft}{\parbox{\dimexpr\linewidth-16pt\relax}{#1}}\endgroup\par}
\newcommand{\WriteBox}[2]{\par\textbf{#1}\par\vspace{2pt}\begin{tikzpicture}\draw[black!40,line width=.5pt] (0,0) rectangle (\dimexpr\linewidth-.5pt\relax,#2);\end{tikzpicture}\par}
\newcommand{\StudentFooter}{\vfill{\small Build first. Explain by pointing, talking, or drawing.}\par}
\newcommand{\Code}[1]{\texttt{#1}}

\hypersetup{pdftitle={Week 6: extra-grades-6-7}}
\begin{document}
\PageTitle{EXTRA / GRADES 6--7}{Five places, only four tests}\NameLine
\RuleBox{Write R as 0 and B as 1. A score still counts correct \textbf{positions}. Choose all tests in advance: \Code{00000}, \Code{00011}, \Code{00101}, \Code{01001}.}
\Task{1. Decode a real example.}{These four tests score 2, 2, 2, 0 in that order. Find the secret. Verify every score.}
\WriteBox{Secret and check:}{.65in}
\Task{2. Why could this always work?}{Let the secret be \(x_1x_2x_3x_4x_5\), with each \(x_i\) equal to 0 or 1. The first score tells the total number of 1s. What does changing two test positions from 0 to 1 tell you?}
\WriteBox{Relate the score changes to these sums: \(x_4+x_5,\ x_3+x_5,\ x_2+x_5\).}{.7in}
\Task{3. The only tricky case.}{A pair-sum of 0 or 2 forces both bits. What if all three pair-sums are 1? Try \(x_5=0\) and \(x_5=1\). Can they have the same total number of 1s, even after you choose \(x_1\)?}
\WriteBox{Why the four answers identify every possible secret:}{.75in}
\textbf{A geometric translation:} Join binary words that differ in one position. Distance is the number of mismatches, so distance = 5 minus score. These four test words are landmarks that distinguish all 32 vertices of a five-dimensional cube.
\StudentFooter
\end{document}
